\documentclass[border=10pt]{standalone}
\usepackage{tikz}
\usepackage{helvet}
\usepackage{xcolor}
\renewcommand{\familydefault}{\sfdefault} % Use sans-serif font

% Load required TikZ libraries
\usetikzlibrary{positioning, shapes.geometric, arrows.meta, calc, shadows}
\usepackage{xcolor}

\definecolor{darkWhite}{rgb}{0.96,0.96,0.96}
\definecolor{bluekeywords}{rgb}{0.13,0.13,1}
\definecolor{greencomments}{rgb}{0,0.5,0}
\definecolor{redstrings}{rgb}{0.9,0,0}

\definecolor{Comment}{RGB}{97,161,176}

\definecolor{btfGreen}{RGB}{51,160,44}
\definecolor{btfRed}{RGB}{190,60,90}

\definecolor{bleuUni}{RGB}{0, 157, 224}
\definecolor{marronUni}{RGB}{68, 58, 49}

\definecolor{bluecite}{HTML}{009DE0}

\definecolor{Paired_1}{RGB}{31,120,180}
\definecolor{Paired_2}{RGB}{166,206,227}
\definecolor{Paired_3}{RGB}{51,160,44}
\definecolor{Paired_4}{RGB}{178,223,138}
\definecolor{Paired_5}{RGB}{227,26,28}
\definecolor{Paired_6}{RGB}{251,154,153}
\definecolor{Paired_7}{RGB}{255,127,0}
\definecolor{Paired_8}{RGB}{253,191,111}
\definecolor{Paired_9}{RGB}{106,61,154}
\definecolor{Paired_10}{RGB}{202,178,214}
\definecolor{Paired_11}{RGB}{177,89,40}
\definecolor{Paired_12}{RGB}{255,255,153}
\definecolor{Accent_1}{RGB}{127,201,127}
\definecolor{Accent_2}{RGB}{190,174,212}
\definecolor{Accent_3}{RGB}{253,192,134}
\definecolor{Accent_4}{RGB}{255,255,153}
\definecolor{Accent_5}{RGB}{56,108,176}
\definecolor{Accent_6}{RGB}{240,2,127}
\definecolor{Accent_7}{RGB}{191,91,23}
\definecolor{Accent_8}{RGB}{102,102,102}
\definecolor{Spectral_1}{RGB}{158,1,66}
\definecolor{Spectral_2}{RGB}{213,62,79}
\definecolor{Spectral_3}{RGB}{244,109,67}
\definecolor{Spectral_4}{RGB}{253,174,97}
\definecolor{Spectral_5}{RGB}{254,224,139}
\definecolor{Spectral_6}{RGB}{255,255,191}
\definecolor{Spectral_7}{RGB}{230,245,152}
\definecolor{Spectral_8}{RGB}{171,221,164}
\definecolor{Spectral_9}{RGB}{102,194,165}
\definecolor{Spectral_10}{RGB}{50,136,189}
\definecolor{Spectral_11}{RGB}{94,79,162}
\definecolor{Set1_1}{RGB}{228,26,28}
\definecolor{Set1_2}{RGB}{55,126,184}
\definecolor{Set1_3}{RGB}{77,175,74}
\definecolor{Set1_4}{RGB}{152,78,163}
\definecolor{Set1_5}{RGB}{255,127,0}
\definecolor{Set1_6}{RGB}{255,255,51}
\definecolor{Set1_7}{RGB}{166,86,40}
\definecolor{Set1_8}{RGB}{247,129,191}
\definecolor{Set1_9}{RGB}{153,153,153}
\definecolor{Set1_10}{RGB}{0,0,0}
\definecolor{Set2_1}{RGB}{102,194,165}
\definecolor{Set2_2}{RGB}{252,141,98}
\definecolor{Set2_3}{RGB}{141,160,203}
\definecolor{Set2_4}{RGB}{231,138,195}
\definecolor{Set2_5}{RGB}{166,216,84}
\definecolor{Set2_6}{RGB}{255,217,47}
\definecolor{Set2_7}{RGB}{229,196,148}
\definecolor{Set2_8}{RGB}{179,179,179}
\definecolor{Dark2_1}{RGB}{27,158,119}
\definecolor{Dark2_2}{RGB}{217,95,2}
\definecolor{Dark2_3}{RGB}{117,112,179}
\definecolor{Dark2_4}{RGB}{231,41,138}
\definecolor{Dark2_5}{RGB}{102,166,30}
\definecolor{Dark2_6}{RGB}{230,171,2}
\definecolor{Dark2_7}{RGB}{166,118,29}
\definecolor{Dark2_8}{RGB}{102,102,102}
\definecolor{Reds_1}{RGB}{255,245,240}
\definecolor{Reds_2}{RGB}{254,224,210}
\definecolor{Reds_3}{RGB}{252,187,161}
\definecolor{Reds_4}{RGB}{252,146,114}
\definecolor{Reds_5}{RGB}{251,106,74}
\definecolor{Reds_6}{RGB}{239,59,44}
\definecolor{Reds_7}{RGB}{203,24,29}
\definecolor{Reds_8}{RGB}{165,15,21}
\definecolor{Reds_9}{RGB}{103,0,13}
\definecolor{Greens_1}{RGB}{247,252,245}
\definecolor{Greens_2}{RGB}{229,245,224}
\definecolor{Greens_3}{RGB}{199,233,192}
\definecolor{Greens_4}{RGB}{161,217,155}
\definecolor{Greens_5}{RGB}{116,196,118}
\definecolor{Greens_6}{RGB}{65,171,93}
\definecolor{Greens_7}{RGB}{35,139,69}
\definecolor{Greens_8}{RGB}{0,109,44}
\definecolor{Greens_9}{RGB}{0,68,27}
\definecolor{Blues_1}{RGB}{247,251,255}
\definecolor{Blues_2}{RGB}{222,235,247}
\definecolor{Blues_3}{RGB}{198,219,239}
\definecolor{Blues_4}{RGB}{158,202,225}
\definecolor{Blues_5}{RGB}{107,174,214}
\definecolor{Blues_6}{RGB}{66,146,198}
\definecolor{Blues_7}{RGB}{33,113,181}
\definecolor{Blues_8}{RGB}{8,81,156}
\definecolor{Blues_9}{RGB}{8,48,107}


\begin{document}

\begin{tikzpicture}[
    % Define consistent styles for the diagram
    base/.style={align=center, font=\large, drop shadow={opacity=0.15, shadow xshift=1mm, shadow yshift=-1mm}, thick},
    % Input nodes (Left)
    % Input nodes (Left) - was ellipse -> rectangle
    input/.style={
        base,
        rectangle,
        rounded corners=2mm,
        draw=black!70,
        fill=white,
        text width=2.8cm,
        minimum height=1.3cm
    },
    % Process nodes (Green, Blue, Orange)
    processG/.style={base, rectangle, rounded corners=2mm, draw=Paired-3, fill=Paired-4, text width=2.4cm, minimum height=2.2cm},
    processB/.style={base, rectangle, rounded corners=2mm, draw=Paired-1, fill=Paired-2, text width=2.6cm, minimum height=2.2cm},
    processO/.style={base, rectangle, rounded corners=2mm, draw=Paired-7, fill=Paired-8, text width=2.4cm, minimum height=2.2cm},
    % Top Configuration nodes - was ellipse/circle -> rectangle
    topInputE/.style={
        base,
        rectangle,
        rounded corners=2mm,
        draw=black!70,
        fill=white,
        text width=2.8cm,
        minimum height=1cm
    },
    topInputC/.style={
        base,
        rectangle,
        rounded corners=2mm,
        draw=black!70,
        fill=white,
        text width=3cm,
        minimum height=1.2cm
    },
    % Output nodes (Right & Bottom) - was ellipse -> rectangle
    output/.style={
        base,
        rectangle,
        rounded corners=2mm,
        draw=Paired-5,
        fill=Paired-6,
        text width=3cm,
        minimum height=1.2cm
    },
    % Connecting Arrows
    arrow/.style={-{Stealth[length=2.5mm, width=2mm]}, thick, draw=black!70}
]

    % --- Place main horizontal processing spine ---
    \node[processG] (export) {ONNX\\Export};
    \node[processB, right=0.8cm of export] (backward) {Backward\\ONNX\\Generation};
    \node[processB, right=0.8cm of backward] (check) {ONNX Check\\Stream Pre-processing};
    \node[processB, right=0.8cm of check] (opt) {Optimizer\\Pass};
    \node[processO, right=1.8cm of opt] (stream) {Stream}; % Extra space for the branch

    % --- Place left initial inputs ---
    \node[input, left=1cm of export, yshift=1.1cm] (nn) {Neural Network\\(CNN,\\Transformers)};
    \node[input, left=1cm of export, yshift=-1.1cm] (batch) {Batch Size,\\Input Size};

    % --- Place top configuration inputs ---
    \node[topInputE, above=1cm of backward] (loss) {Loss, Frozen\\Weights};
    \node[topInputC, above=1cm of opt] (opt_cfg) {Optimizer\\configuration\\(Adam, SGD...)};
    \node[topInputC, above=1cm of stream] (accel_cfg) {Accelerator\\Mapping\\Layer-Fused\\Configurations};

    % --- Place bottom/right outputs ---
    % Calculate midpoint for the branching output
    \coordinate (branch_pt) at ($(opt.east)!0.45!(stream.west)$);
    \node[output, below=1.5cm of branch_pt] (train) {Training\\ONNX};
    
    \node[output, right=1.2cm of stream, yshift=1.1cm] (latency) {Latency,\\Energy,\\Memory};
    \node[output, right=1.2cm of stream, yshift=-1.1cm] (schedule) {Schedule};

    % --- Draw Connections ---
    % Left Inputs to ONNX Export
    \draw[arrow] (nn.east) -- ([yshift=0.5cm]export.west);
    \draw[arrow] (batch.east) -- ([yshift=-0.5cm]export.west);

    % Main horizontal spine
    \draw[arrow] (export.east) -- (backward.west);
    \draw[arrow] (backward.east) -- (check.west);
    \draw[arrow] (check.east) -- (opt.west);
    
    % Draw arrow from Optimizer to Stream, and create a branch point
    \draw[arrow] (opt.east) -- (stream.west);
    
    % Top inputs downwards
    \draw[arrow] (loss.south) -- (backward.north);
    \draw[arrow] (opt_cfg.south) -- (opt.north);
    \draw[arrow] (accel_cfg.south) -- (stream.north);

    % Branching to Training ONNX
    % Draws a line from the calculated branch point straight down to the node
    \draw[arrow] (branch_pt) -- (train.north);
    % To make the T-junction look solid, draw a small line segment
    \filldraw[black!70] (branch_pt) circle (1.5pt);

    % Right Outputs
    \draw[arrow] (stream.east) -- (latency.west);
    \draw[arrow] (stream.east) -- (schedule.west);

\end{tikzpicture}

\end{document}